\documentclass[11pt]{amsart}

\usepackage[T1]{fontenc}
\usepackage{lmodern}
\usepackage{microtype}
\usepackage{amsmath,amssymb,amsthm}
\usepackage[hidelinks]{hyperref}

\hypersetup{
  pdftitle={A Counterexample to the Jahanbekam--West Matching Conjecture}
}

\newtheorem{theorem}{Theorem}
\theoremstyle{remark}
\newtheorem*{remark}{Remark}

\title[Counterexample to the Jahanbekam--West conjecture]
{A Counterexample to the Jahanbekam--West Matching Conjecture}

\subjclass[2020]{Primary 05C70; Secondary 05C15}
\keywords{rainbow perfect matching, edge coloring, Jahanbekam--West conjecture, counterexample}
\date{}

\begin{document}
\raggedbottom

\begin{abstract}
Let $m(n,t)$ be the maximum number of colors in an edge-coloring of $K_n$
containing no $t$ edge-disjoint rainbow perfect matchings. Jahanbekam and
West conjectured that
\[
  m(n,t)=\binom{n-2}{2}+t
\]
for even $n>4$ and $2\le t<n$. For every even $n\ge6$, we construct a
coloring of $K_n$ with $\binom{n-2}{2}+n$ colors and no $n-1$
edge-disjoint rainbow perfect matchings. Thus the conjecture fails for
$t=n-1$.
\end{abstract}

\maketitle

A subgraph of an edge-colored graph is \emph{rainbow} if its edges have
pairwise distinct colors. For even $n$ and $2\le t<n$, let $m(n,t)$ denote
the maximum number of colors in an edge-coloring of $K_n$ containing no
$t$ edge-disjoint rainbow perfect matchings. Jahanbekam and West
\cite[Theorem~1.1]{JW} proved that
\[
  m(n,t)\ge \binom{n-2}{2}+t,
\]
with equality when $n>4t+10$. Immediately after that theorem they write
``We believe that equality holds whenever $n>t\ge2$ if $n>4$; note that
$m(n,t)=\binom{n-2}{2}+t+1$ when $n=4$'', and at the start of
\cite[Section~4]{JW}, having described the construction that gives the
lower bound, ``We conjecture that the construction is optimal for
$n>t$''. The conjectured range thus includes $t=n-1$, whereas the range
$n>4t+10$ in which equality is proved never does, since $t=n-1$ there
would require $n>4n+6$.

The exception at $n=4$ recorded in the first of these sentences is the
seed of the coloring below. For $n=4$ the set $B$ defined in the proof is
itself a perfect matching of $K_4$; giving its two edges a common color
leaves $1+\binom{4}{2}-2=5=\binom{2}{2}+3+1$ colors and, since the three
perfect matchings of $K_4$ partition $E(K_4)$, no three edge-disjoint
rainbow perfect matchings. That coloring therefore realizes the value $m(4,3)$
recorded by Jahanbekam and West, and what follows is the observation that
the same configuration keeps working at $t=n-1$ for every even $n\ge6$.
Accordingly, the conjecture should be restated as
$m(n,t)=\binom{n-2}{2}+t$ for even $n>4$ and $2\le t\le n-2$.

\begin{theorem}
For every even integer $n\ge6$,
\[
  m(n,n-1)\ge \binom{n-2}{2}+n
  >\binom{n-2}{2}+(n-1).
\]
Consequently, the Jahanbekam--West matching conjecture is false at
$t=n-1$.
\end{theorem}

\begin{proof}
Choose distinct vertices $x,y,z\in V(K_n)$ and set
\[
  B=\{xv:v\in V(K_n)\setminus\{x,y,z\}\}\cup\{yz\}.
\]
Color all edges in $B$ with one color, and give each edge outside $B$ a
distinct new color. Since $|B|=n-2$, this coloring uses
\[
  1+\binom{n}{2}-(n-2)=\binom{n-2}{2}+n
\]
colors.

Suppose that $M_1,\ldots,M_{n-1}$ are pairwise edge-disjoint perfect
matchings. They contain
\[
  (n-1)\frac n2=\binom n2
\]
edges altogether, and hence partition $E(K_n)$. Let $M_j$ be the matching
containing $yz$. Since $y$ and $z$ are matched to each other in $M_j$,
the vertex $x$ is matched in $M_j$ to some
$v\in V(K_n)\setminus\{x,y,z\}$. Thus $M_j$ contains both $yz$ and $xv$,
which have the same color. Therefore $M_j$ is not rainbow. The coloring
contains no $n-1$ edge-disjoint rainbow perfect matchings.
\end{proof}

\begin{remark}
The construction gives only the displayed lower bound for $m(n,n-1)$; it
does not determine this parameter exactly. Nor does it bear on $t\le n-2$.
Indeed, delete from a $1$-factorization of $K_n$, which exists because $n$
is even, the factor containing $yz$. Each of the remaining $n-2$ factors
avoids $yz$ and contains exactly one edge at $x$, so it meets $B$ in at
most one edge and is therefore rainbow. Hence the coloring above does
admit $n-2$ pairwise edge-disjoint rainbow perfect matchings, and the
conjecture is untouched for $t\le n-2$.
\end{remark}

\begin{thebibliography}{1}

\bibitem{JW}
S.~Jahanbekam and D.~B. West,
\newblock \emph{Anti-Ramsey problems for $t$ edge-disjoint rainbow spanning
subgraphs: cycles, matchings, or trees},
\newblock J. Graph Theory \textbf{82} (2016), no.~1, 75--89,
\newblock
\href{https://doi.org/10.1002/jgt.21888}{doi:10.1002/jgt.21888}.
\newblock Quoted here from the authors' revised version of April 2015,
\url{https://dwest.web.illinois.edu/pubs/aramcmt.pdf}.

\end{thebibliography}

\end{document}